Different from the creative generation, grounded article generation may impact how people learn about topics or consume source information. All the studies and the evaluation in this work are designed to prevent the dissemination of misinformation by not publishing generated content online and implementing strict accuracy checks. We avoid any disruption to Wikipedia or related communities, as our system does not interact with live pages. Also, although we try to generate grounded articles, we believe there is no privacy issue related to this work as we only use information publicly available on the Internet.

The primary risk of our work is that the Wikipedia articles written by our system are grounded on information on the Internet which contains some biased or discriminative content on its own. Currently, our system relies on the search engine to retrieve information but does not include any post-processing module. We believe improving the retrieval module to have good coverage of different viewpoints and adding a content sifting module to the current system will be a critical next step to achieve better neutrality and balance in the generated articles. %

Another limitation we see from an ethical point of view is that we only consider writing English Wikipedia articles in this work. Extending the current system to a multilingual setup is a meaningful direction for future work as more topics do not have Wikipedia pages in non-English languages.
